\subsection{Классификация.}
\begin{itemize}
	\item ГУ 1\textsuperscript{го} типа(рода): $ u|_{x=0}=\mu(t),\,t\ge0 $ -- заданное решение
	\item ГУ 2\textsuperscript{го} типа(рода): $ u_x|_{x=0}=\mu(t),\,t\ge0 $ -- заданная сила \\
	      \hspace*{2em}$\displaystyle T_0 u_x(0,t)=-\Phi_0(t) \;\Longrightarrow\; \nu(t)=-\frac{\Phi_0(t)}{T_0} $ %FIXME кака 
	\item Однородное ГУ 2\textsuperscript{го} типа: $ u_x|_{x=0}=0,\,t\ge0 $ -- свободный конец
	\item Однородное ГУ 1\textsuperscript{го} типа: $ u|_{x=0}=0$ -- зарепленный конец%FIXME а ограничения на t есть?
	\item ГУ 3\textsuperscript{го} типа(рода): $ (u+\alpha u_x)|_{x=0}=\phi(t),\,t\ge0 $ -- закрепление+сила%FIXME управление??
	\item Однородное ГУ 3\textsuperscript{го} типа(рода): $ (u+\alpha u_x)|_{x=0}=0$ -- закрепление+сила%FIXME управление??
\end{itemize}
\subsection{Задача Коши для уравнения колебания бесконечной струны.}
\begin{equation}
	\frac{\partial^2 u}{\partial t^2}=a^2\frac{\partial^2 u}{\partial t^2}+f(x,t),\qquad x\in\mathbb{R},t>0
	\tag{$1,f$}
\end{equation}
\begin{equation}
	\text{НУ:}\qquad u|_{t=0}= u_0(x),\quad u_t|_{t=0}=u_1(x),\quad x \in \mathbb{R}
	\tag{2}
\end{equation}
$$ u=u(x,t)\in C^2(t>0) \land C^1(t\ge0),$$
$$f(x,t)\in C^1(t\ge0),\quad u_0(x)\in C^2(\mathbb R),\quad u_1(x)\in C^1(\mathbb R) $$
\begin{theorem}[Существование и единственность]
	Задача Коши для УКБС $(1.f),\,(2)$ в рассматриваемом классе функций имеет решение, притом единственное, которое представдяется
	формулой \emph{Даламбера}(Ф.Д.f):
	$$
		u(x,t)=\frac{u_0(x-at)+u_0(x+at)}2 +\frac{1}{2a} \int^{x+at}_{x-at}u_1(\xi)\mathrm{d}\xi+
		\frac{1}{2a}\int^t_0\mathrm d\tau \int^{x+a(t-\tau)}_{x+a(t-\tau)}f(\xi,\tau)\mathrm d\xi
	$$
\end{theorem}
Рассмотрим задачу коши для ОУКБС:
\begin{equation}
	u_{tt}= au_{xx},\quad x \in \mathbb{R},\;t>0
	\tag{1.o}
\end{equation}
$(2)$, то есть:
$$Lu=u_{tt}-a^2u_{xx}$$
Уравнение характеристик:
\begin{equation}
	\mathrm d x^2-a^2\mathrm dt^2=0\rightarrow (\mathrm d x-a\mathrm d t)(\mathrm d x+a\mathrm d t)\rightarrow
	\left[ \begin{array}{ll}
		C_1=x-at \\
		C_2=x+at
	\end{array} \right.
	\tag*{--- первые интеграллы}
\end{equation}
Переходим к характерестическим переменным $\displaystyle\xi,\eta:\begin{array}{ll}\xi=x-ax\\\eta=x+at\end{array}\rightarrow$
приводится к виду: $\tilde u_{\xi\eta}=0$\\%FIXME чет с русским
его решение: $\displaystyle \tilde u(\xi,\eta)=\tilde f(\xi)+\tilde g(\eta)$ или $u(x,t)=f(x-at)+g(x+at)$ -- ОРУ,
где $\tilde f$ и $\tilde g$ -- достаточно гладкие функции переменных $\xi$ и $\eta$.\\
Решение задачи Коши найдем, подобрав соответствующие функции $\tilde f(\xi)$ и $\tilde g(\eta)$ так, чтобы удовлетворять НУ $(2)$.\\
найдем $u_t$:\\
$u_t(x,t)=-af'(x-at)+ag'(x+at)$, тогда из НУ:\\
$\displaystyle
	\left\{\begin{array}{ll}
		u|_{t=0}=f(x)+g(x)=u_0(x),        & x\in\mathbb R \\
		u_t|_{t=0}=-af'(x)+ag'(x)=u_1(x), & x\in\mathbb R
	\end{array}\right. \Longrightarrow -f(x)+g(x)=\frac1a \int^x_0 u_1(\xi)\mathrm d\xi+C,\,x\in\mathbb R
$\\
Таким образом находим:
$$f(x)=\frac12u_0(x)-\frac1{2a}\int^x_0 u_1(\xi)\mathrm d\xi-\frac C2,\,x \in\mathbb R$$
$$g(x)=\frac12u_0(x)-\frac1{2a}\int^x_0 u_1(\xi)\mathrm d\xi+\frac C2,\,x \in\mathbb R$$
cледовательно решение задачи коши для ОУКБС имеет вид:
\begin{equation}
	u(x,t)=\frac{u_0(x-at)+u_0(x+at)}2 +\frac{1}{2a} \int^{x+at}_{x-at}u_1(\xi)\mathrm{d}\xi,\quad x\in\mathbb R,\,t\ge0
	\tag*{--- формула Даламбера(однородная)}
\end{equation}
\underline{\emph{Физический смысл}}:  ``метод распределения волн''.\\%FIXME КАРТИИИИИИИИИИИИНКААААААААА
Рассмотрим ОРУКБС вида
$$u(x,t)=f(x-at),\;g=0$$
Пусть наблюдатель движется в положительном направлении со скоростью $a$ и $x(t=0)=0$.
Свяжем с наблюдателем $O'x't'\;\Longrightarrow\; x=x'+at',\,t=t'$ (\emph{преобразование Галилея}) или $x'=x-at,\,t'=t\;\Longrightarrow$ в сиситеме координат $O'x't'$, связанной с наблюдателем решение будет иметь вид:
$u=f(x')$, то есть наблюдатель движется в право со скоростью $a$. В его системе координат будет наблюдаться постоянный профиль струны, следовательно функция $f(x-at)$ описывает волну, бегущую вправо со скоростью $a$.\\
Аналогично, функция $g(x+at)$ \textbf{---//---}, бегущую влево со скоростью $a$.\\
Таким образом, в соответствии с формулой Даламбера(однородной):
$$ f,g(s)=\frac12u_0(s)\pm\frac1{2a}\int_0^xu_1(\xi)\mathrm d\xi\pm\frac C2 $$
\begin{note}
	Пусть решение ОУКБС $u(x,t)=f(x-at)$, причем $f(x)=0$ вне интервала $(x_1,\ x_2)$, то есть $f(x)$ -- финитная функция.\\
	Тогда характеристики $x-at=x_1$ и $x-at=x_2$ разбивают $Oxt$ на 3 подобласти,%FIXME КРТИНКА
	при это $u(x,t)=0$ в I и III подоблостях, а характеристики $x-at=x_1$ и $x-at=x_2$ зазывают \emph{задним} и \emph{передним фронтами волны}.
\end{note}
Рассмотрим формулу Даламбера для однородного уравнения. Зафиксируем $M_0(x_0,t_0)$ и проведем через нее характеристики до пересечения с осью $x$. Получим точки:\\
$\displaystyle\begin{array}{ll}P:\,x_1=x_0-at_0,\;t_1=0\\Q:\,x_2=x_0+at_0,\;t_2=0\end{array}$, с характеристиками
$\displaystyle\begin{array}{ll}x+at=x_0-at_0\\x+at=x_0+at_0\end{array}$%КАРТИИИНКА из лекций V2
\begin{definition}[Характерестический треугольник и область зависимости]
	Треугольник $MPQ$ называют \emph{характерестическим треугольником} для точки $M$, а интервал $(x_1,x_2)$ -- \emph{областью зависимости} точки $M_0$.
\end{definition}
При этом, значение отклонение струны $u(x_0.t_0)$ точки $M_0$ определяют лишь начальные условия на отрезке $[x_1,x_2]$.\\
\textbf{\underline{Вывод:}} если Н.У. заданы лишь на отрезке $[x_1,x_2]$, то решение ЗК для ОУКБС однозначно определяется в характерестическом треугольник $MPQ$, в котором $[x_1,x_2]$ является областью зависимости, а соторы ограничены характеристиками.

